\providecommand{\BM}{\textbf{[B]}}
\providecommand{\KM}{\textbf{[K]}}
\providecommand{\SM}{\textbf{[S]}}
\providecommand{\XM}{\textbf{[X]}}
\providecommand{\QM}{\textbf{[Q]}}

\chapter*{The Hierarchy $v/E_P$: Two Derivations Compared}
\addcontentsline{toc}{chapter}{The Hierarchy $v/E_P$: Two Derivations Compared}

{\small\noindent\textbf{Abstract.}\quad\ignorespaces
FFGFT derives the electroweak scale from $\xi$ (R104) and reaches the Planck
scale through its own chain $\xi\to G\to E_P$. FFGFT therefore delivers the
hierarchy number $v/E_P$ entirely by itself; the electron mass serves only as
the SI anchor. Observer Patch Holography (OPH) delivers the same number
conditionally through its Theorem~A. The two derivations give
$2.0389\times10^{-17}$ (FFGFT, $+1.10\,\%$) and $2.0200\times10^{-17}$ (OPH,
$+0.16\,\%$). The FFGFT deviation is entirely the known bare gap of the lepton
ladder (Doc.~352); the Planck chain contributes less than $10^{-4}$. The FFGFT
entry for a hashed comparison table is given.
}\medskip

\section*{List of symbols}
\addcontentsline{toc}{section}{List of symbols}

\medskip
\begin{tabular}{@{}ll@{}}
\textbf{[B]} & algebraically proven \\
\textbf{[K]} & numerically confirmed \\
\textbf{[S]} & structurally plausible, derivation open \\
\textbf{[Q]} & taken from an external source, cited \\
\end{tabular}

\medskip
\begin{tabular}{@{}lp{10cm}@{}}
\toprule
\textbf{Symbol} & \textbf{Meaning} \\
\midrule
$\xi$ & FFGFT scaling parameter ($4/30000$) \\
$v$ & electroweak scale \\
$E_P$ & Planck energy (not reduced) \\
$G$, $C_{\rm conv}$, $K_{\rm frak}$ & gravitational constant and factors of the SI representation (Docs.~012/013) \\
$E_\star$, $P$, $\alpha_U$ & OPH reference energy, pixel ratio, unified coupling \\
\bottomrule
\end{tabular}

\section{Occasion}

Doc.~374 brought FFGFT and OPH together at the electroweak scale and, in doing
so, combined $m_e/v$ from the FFGFT ladder with $v/E_\star$ from OPH. That
understates the FFGFT side: FFGFT derives $v$ itself (R104) and reaches $E_P$
through its own chain (Doc.~374, §\,5). The hierarchy number $v/E_P$ is
therefore not a combination but a second, independent derivation of the same
quantity. This document places the two derivations side by side.

\section{The FFGFT derivation}

\begin{enumerate}\itemsep2pt
\item Electroweak scale (R104, Doc.~006):
  \begin{equation}
    v=\frac{m_e}{\tfrac43\,\xi^{3/2}}=248.93\,\mathrm{GeV}.
  \end{equation}
\item Planck energy through its own chain (Docs.~012, 013, 180):
  $G=\xi^2/(4m_e)\cdot C_{\rm conv}\cdot K_{\rm frak}$, hence
  $E_P=1.22087\times10^{19}$\,GeV.
\item Hierarchy:
  \begin{equation}
    \left.\frac{v}{E_P}\right|_{\rm FFGFT}=2.0389\times10^{-17}\quad\KM .
  \end{equation}
\end{enumerate}

Against the SM value $v=246.22$\,GeV (definition from $G_F$, R112) this is
$+1.10\,\%$ higher. The deviation is entirely that of $v$; the Planck chain
contributes less than $10^{-4}$. It is identical to the bare gap of the
lepton ladder $m_e/v$ ($-1.09\,\%$, Doc.~352): $v$ is obtained from the bare
relation and inherits its residual. Not a new finding.

\section{The OPH derivation \QM}

Theorem~A from \texttt{code/particles/hierarchy/theorem\_package.md}:
\begin{equation}
  \frac{v}{E_\star}=P^{-1/2}\exp\!\left[-\frac{2\pi}{4\alpha_U(P)}\right]
  =2.0200\times10^{-17},
\end{equation}
reproduced in the check script. The derivation is explicitly conditional in
OPH; the identification $E_\star=E_P$ is a reading \SM{} (Doc.~374). Against
the SM value: $+0.16\,\%$.

\section{Comparison}

\begin{center}
\begin{tabular}{@{}lrrl@{}}
\toprule
\textbf{Derivation} & $v/E_P$ & \textbf{vs.\ SM} & \textbf{Origin of the deviation} \\
\midrule
FFGFT & $2.0389\times10^{-17}$ & $+1.10\,\%$ & bare residual of the ladder (Doc.~352) \\
OPH \QM & $2.0200\times10^{-17}$ & $+0.16\,\%$ & conditional fixed point, reading $E_\star=E_P$ \\
SM comparator & $2.0167\times10^{-17}$ & --- & $v$ from $G_F$ (R112) \\
\bottomrule
\end{tabular}
\end{center}

The two derivations lie $0.93\,\%$ apart. They also differ in sensitivity:
\begin{itemize}\itemsep2pt
\item FFGFT: $\partial\ln v/\partial\ln\xi=-3/2$ \BM. A rational structural
  quantity with power $3/2$; small changes of $\xi$ act weakly.
\item OPH: $\partial\ln(v/E_\star)/\partial\alpha_U=\pi/(2\alpha_U^2)\approx929$.
  The coupling sits in the exponent; the $0.16\,\%$ corresponds to a shift of
  $\alpha_U$ by $4\times10^{-5}$ relatively.
\end{itemize}
The FFGFT value is robust and carries a known residual; the OPH value is more
accurate but highly sensitive and conditional. Both statements are testable:
on the FFGFT side through the resolution of the bare residual (Doc.~352), on
the OPH side through the clock attachment and the identification of
$E_\star$.

\section{Entry for the comparison table}

For a hashed comparison table the FFGFT entry, including the route, reads:

\begin{center}
\begin{minipage}{0.92\textwidth}
\small\ttfamily
v/E\_P | FFGFT | 2.03894e-17 | route: v=m\_e/((4/3)xi\^{}(3/2)) (R104, Doc 006) ; E\_P via G=xi\^{}2/(4m\_e)C\_conv K\_frak (Docs 012/013/180) ; bare ladder
\end{minipage}
\end{center}

\noindent SHA-256 of the string (generated in the check script):\\
{\small\ttfamily 84fac85fc55d8c5928a07550e774683961613059415b0ef02715b327f58fb15f}

\section*{Tool and method}
\addcontentsline{toc}{section}{Tool and method}
Claude (Anthropic) was used as a tool for: writing and running the Python
check script, drafting the \LaTeX{} sources, and maintaining the changelog.
The tool computes and writes; the author decides what is computed and
written, judges every result and determines what it means physically. OPH
values are taken verbatim from the public repository; measured values
(CODATA 2022, PDG) serve only as comparators.

\section{Result}

\begin{enumerate}
\item FFGFT delivers $v/E_P=2.0389\times10^{-17}$ entirely by itself: $v$ from
  $\xi$ (R104), $E_P$ from its own chain \KM.
\item The deviation of $+1.10\,\%$ is the bare residual of the lepton ladder;
  the Planck chain contributes less than $10^{-4}$ \KM.
\item OPH delivers, conditionally, $2.0200\times10^{-17}$ ($+0.16\,\%$) \QM;
  the derivations lie $0.93\,\%$ apart.
\item Sensitivity: FFGFT $-3/2$ in $\ln\xi$ \BM, OPH $\approx929$ in $\alpha_U$.
\item FFGFT table entry given with route and SHA-256.
\end{enumerate}

\medskip
\noindent\textbf{Check script:} \texttt{2/python/Dok375\_Skripte/pruef\_375\_hierarchie.py} (11/11)

\medskip
\noindent\textbf{Register (Doc.~190):} no entry.

\begin{thebibliography}{9}
\bibitem{dok006} J.~Pascher, \emph{Doc.~006 --- T0 particle masses} (2025).
\bibitem{dok012} J.~Pascher, \emph{Doc.~012 --- T0 gravitational constant} (2025).
\bibitem{dok013} J.~Pascher, \emph{Doc.~013 --- T0 and SI units} (2025).
\bibitem{dok180} J.~Pascher, \emph{Doc.~180 --- Derivation of $L_0$} (2026).
\bibitem{dok352} J.~Pascher, \emph{Doc.~352 --- Lepton residuals} (2026).
\bibitem{dok374} J.~Pascher, \emph{Doc.~374 --- Common anchor} (2026).
\bibitem{oph} B.~Mueller et al., \emph{Observer Patch Holography},
  \texttt{github.com/FloatingPragma/observer-patch-holography} (as of 22~September~2026).
\end{thebibliography}
